\documentclass[11pt]{article}
\usepackage[margin=1in]{geometry}
\usepackage{amsmath,amssymb,amsthm,mathtools}
\usepackage[hidelinks]{hyperref}
\newtheorem{theorem}{Theorem}
\newtheorem{lemma}[theorem]{Lemma}
\newtheorem{corollary}[theorem]{Corollary}
\newcommand{\spen}{s^{+}}
\title{Positive Square Energy of Theta Graphs with One Tree Attachment}
\author{math-god}
\date{24 July 2026}
\begin{document}
\maketitle

\begin{abstract}
We prove that if $G$ is a connected graph whose 2-core is a simple theta
graph $\Theta(a,b,c)$ and one arbitrary rooted tree is attached to a single
core vertex through one bridge edge, then $s^+(G)>|V(G)|$.
This extends the bare theta theorem of our companion paper to the
weighted-core setting of Conjecture~1.2 of Akbari et al.\ (arXiv:2506.07264).
The proof combines an exact root-congruence PSD witness, a local
four-fifths reduction via polynomial majorants, and a complete phase-area
classification of the signed square trace by odd-girth residue class.
\end{abstract}

\section{Introduction}
For a graph $G$ with adjacency eigenvalues $\lambda_i$, let
$s^+(G)=\sum_{\lambda_i>0}\lambda_i^2$.
The companion paper proves that every simple theta graph $G$ satisfies
$s^+(G)>|V(G)|$.
Here we extend this to one arbitrary rooted-tree attachment.

\begin{theorem}\label{thm:weighted}
Let $H$ be a simple theta graph and let $v\in V(H)$.
Let $T$ be a rooted tree with root $r$ of degree one in $T-r$,
and let $G$ be obtained by identifying $r$ with $v$ and adding one
bridge edge from $r$ to a new vertex.
Then
\[ s^+(G)>|V(G)|. \]
More precisely, for every root $v$ of $H$,
\[ s^+\!\left(A(H)-\tfrac12 E_{vv}\right)>|V(H)|. \]
\end{theorem}

The weighted 2-core reduction of Akbari et al.\ reduces the attachment
problem to the endpoint $d=1/2$ because $s^+(A-dE_{vv})$ is nonincreasing
in $d$ and the rooted-tree penalty $d=(A^-(T))_{vv}\le 1/2$.

\section{The root-congruence witness}
Let $A$ be the adjacency matrix of $H$, $P=A_+$, and fix a root $v$.
Put $S=\operatorname{tr}(P^2)=s^+(H)$, $a=P_{vv}$, $q=(P^2)_{vv}$.
For $0\le k\le 1$, set $C_k=I-(1-k)E_{vv}$ and $X_k=C_kPC_k\succeq0$.

\begin{lemma}\label{lem:witness}
\[ s^+\!\left(A-\tfrac12E_{vv}\right)\ge
S-\frac{2(1-k)^2q}{1}-k^2a-(1-k^2)^2a^2. \]
At $k=6/7$, this gives
\[ s^+\!\left(A-\tfrac12E_{vv}\right)\ge S-L, \quad
L=\frac{2q}{49}+\frac{36a}{49}+\frac{169a^2}{2401}. \]
\end{lemma}

The proof is direct trace expansion using $AP=P^2$ and $A_{vv}=0$.
The exact certificate is in \texttt{root\_congruence\_certificate.py}.

\section{Local four-fifths reduction}

\begin{lemma}\label{lem:local}
For every rooted simple theta, $L<4/5$.
More precisely, $4/5-L\ge 31684811/1355316480>0$.
\end{lemma}

\begin{proof}
The quadratic $V_t(x)=t/4+x/2+x^2/(4t)\ge x_+$ for all $x$, giving
$a\le t/4+\deg(v)/(4t)$.
The cubic $U(x)=1/5+x/2+x^2/2+x^3/9\ge x_+^2$ on $[-3,3]$,
giving $q\le 1/5+\deg(v)/2+\tau/9$, where $\tau=(A^3)_{vv}=2T_v$.
The five possible $(\deg(v),\tau)$ states are
$(3,0),(3,2),(3,4),(2,0),(2,2)$.
Substituting the sharp bounds gives strict rational margins below $4/5$
in all five cases.
\end{proof}

The exact certificate is in \texttt{local\_four\_fifths\_certificate.py}.
Together, Lemmas~\ref{lem:witness}--\ref{lem:local} prove:

\begin{corollary}\label{cor:conditional}
If $s^+(H)-|V(H)|\ge 4/5$, then
$s^+(A(H)-E_{vv}/2)>|V(H)|$ for every root $v$.
\end{corollary}

\section{The bare four-fifths theorem}

\begin{theorem}\label{thm:bare}
Every simple theta graph $H$ satisfies $s^+(H)-|V(H)|\ge 4/5$,
except exactly $\Theta(2,2,3)$, $\Theta(2,3,3)$, and $\Theta(1,4,4)$.
\end{theorem}

The proof uses the imaginary-axis phase representation.
For $z=e^{-2u}\in(0,1)$ and $x=2i\sinh u$, the theta characteristic
polynomial factors through a phase carrier
\[ H(z)=R(z)+i\sqrt z\,S(z), \]
where $R\ge0$ (in fact $R>0$ for nonbipartite thetas) and
\[ \operatorname{tr}(A|A|)=-\frac2\pi\int_0^1(z^{-2}-1)\operatorname{Arg} H(z)\,dz. \]
Since $s^+-n=1+\operatorname{tr}(A|A|)/2$, the target $4/5$ is equivalent to
\[ \int_0^1(z^{-2}-1)\alpha(z)\,dz\le\frac\pi5. \]

\subsection*{Phase sign classification}
The shortest odd and even path lengths $o,e$ determine
$\operatorname{sign}(PQ)=(-1)^{(o+e-1)/2}$.
If the shortest odd cycle $g=o+e\equiv3\pmod4$, then
$\operatorname{tr}(A|A|)>0$ and $s^+-n>1$.
Bipartite thetas have $s^+-n=1$.
Thus only $g\equiv1\pmod4$ requires quantitative control.

\subsection*{The $\tilde C_5$ class}
When $g=5$, the only residue pairs are $(1,4)$ and $(2,3)$.
Exact mod-four phase monotonicity (\texttt{theta\_phase\_monotonicity\_certificate.py}),
safe common limits ($I_\infty<8/15<\pi/5$), and exact Sturm gates
(\texttt{theta\_short\_base\_gates.py}) settle both channels
outside $\Theta(2,2,3)$, $\Theta(2,3,3)$, $\Theta(1,4,4)$.

\subsection*{Odd girth at least nine}
When $g\ge9$, a 16-case parameterized certificate
(\texttt{theta\_g9\_signed\_certificate.py}) proves
$K=4z^4R-S\ge0$ on $[0,3/4]\times[0,1]^3$ by exact tensor Bernstein
coefficients, giving $\alpha(z)<4z^{9/2}$ on $[0,3/4]$.
On $[3/4,1]$, $\alpha\le\pi/2$ and the tail weight has mass $1/12$.
Thus $I<3051\sqrt3/19712+\pi/24<\pi/5$.

\section{Exceptional weighted endpoints}
For the three bare exceptions, direct Sturm isolation of
$\chi(A-E_{vv}/2)$ at every root orbit proves
$s^+(A-E_{vv}/2)>|V(H)|$ with exact rational margins.
The complete certificate covering all 9 root orbits is in
\texttt{theta\_short\_base\_gates.py}.

\section{Conclusion}
Theorem~\ref{thm:bare} plus Corollary~\ref{cor:conditional} prove
Theorem~\ref{thm:weighted} for every nonexceptional theta.
The three exceptions are covered by their direct weighted certificates.
This completes the one-tree weighted-theta extension.

\begin{thebibliography}{9}
\bibitem{AKMPZ}
S.~Akbari, H.~Kumar, B.~Mohar, S.~Pragada, and S.~Zhang,
\emph{Refinement of a conjecture on positive square energy of graphs},
arXiv:2506.07264v1, 2025.
\end{thebibliography}
\end{document}
